% ===== PRÉAMBULE COMMUN — À COPIER TEL QUEL EN TÊTE DE CHAQUE FICHIER .tex =====
% Ne modifier QUE :
%   - les deux lignes \fancyhead[L] et \fancyhead[R]
%   - le bloc titre \begin{center}...\end{center} juste après \begin{document}
\documentclass[11pt,a4paper]{article}
\usepackage[utf8]{inputenc}
\usepackage[T1]{fontenc}
\usepackage{lmodern}
\usepackage[french]{babel}
\usepackage{amsmath,amssymb,mathtools}
\usepackage{icomma}
\usepackage{xcolor}
\usepackage{tikz}
\usepackage{pgfplots}
\usepackage[most]{tcolorbox}
\usepackage{enumitem}
\usepackage{array,booktabs,colortbl}
\usepackage[a4paper,margin=2.2cm,headheight=26pt,headsep=10pt]{geometry}
\usepackage{fancyhdr}
\usepackage[hidelinks]{hyperref}
\usetikzlibrary{arrows.meta,calc,angles,quotes,positioning,decorations.markings,intersections,backgrounds}
\pgfplotsset{compat=1.18}

\definecolor{primary}{RGB}{226,74,88}
\definecolor{primarydark}{RGB}{150,28,42}
\definecolor{methcol}{RGB}{40,90,150}
\definecolor{excol}{RGB}{0,135,90}
\definecolor{defcol}{RGB}{0,114,178}
\definecolor{attncol}{RGB}{200,40,55}
\definecolor{thmcol}{RGB}{0,140,120}
\definecolor{courbe}{RGB}{32,110,200}
\definecolor{courbeder}{RGB}{210,70,40}
\definecolor{tang}{RGB}{0,150,90}
\definecolor{grille}{RGB}{200,205,215}
\definecolor{plus}{RGB}{200,60,60}
\definecolor{moins}{RGB}{30,90,200}

\tcbset{boxbase/.style={enhanced,breakable,boxrule=0.9pt,arc=2.5pt,
  left=3mm,right=3mm,top=2.3mm,bottom=2.3mm,fonttitle=\bfseries,coltitle=white}}
\newtcolorbox{cle}[1][]{boxbase,colback=defcol!6,colframe=defcol,
  title={\textsc{À retenir}\ifx\relax#1\relax\relax\else\textmd{ : \itshape #1}\fi}}
\newtcolorbox{methode}[1][]{boxbase,colback=methcol!7,colframe=methcol,sharp corners,
  title={\textsc{Méthode}\ifx\relax#1\relax\relax\else\textmd{ --- \itshape #1}\fi}}
\newtcolorbox{exemplebox}[1][]{boxbase,colback=excol!5,colframe=excol,
  title={\textsc{Exemple}\ifx\relax#1\relax\relax\else\textmd{ : \itshape #1}\fi}}
\newtcolorbox{piege}[1][]{boxbase,colback=attncol!6,colframe=attncol,
  title={\textsc{Attention}\ifx\relax#1\relax\relax\else\textmd{ : \itshape #1}\fi}}
\newtcolorbox{exo}[1][]{boxbase,colback=primary!4,colframe=primary,
  title={\textsc{Exercice}\ifx\relax#1\relax\relax\else\textmd{~#1}\fi}}
\newtcolorbox{corr}[1][]{boxbase,colback=thmcol!5,colframe=thmcol,
  title={\textsc{Corrigé}\ifx\relax#1\relax\relax\else\textmd{~#1}\fi}}

\newcommand{\R}{\mathbb{R}}
\newcommand{\N}{\mathbb{N}}
\newcommand{\ds}{\displaystyle}
\newcommand{\tg}{\operatorname{tg}}
\newcommand{\cotg}{\operatorname{cotg}}
\newcommand{\arctg}{\operatorname{arctg}}
\newcommand{\arccotg}{\operatorname{arccotg}}
\newcommand{\limh}{\lim_{h\to 0}}

\pagestyle{fancy}\fancyhf{}
\fancyhead[L]{\small\textcolor{primarydark}{\textbf{Thème 1} — Nombre dérivé \& taux d'accroissement}}
\fancyhead[R]{\small\textcolor{primarydark}{\textsc{Corrigé — Difficile}}}
\fancyfoot[C]{\small\thepage}
\renewcommand{\headrulewidth}{0.8pt}

\begin{document}

\begin{center}
{\LARGE\bfseries\color{primarydark} Corrigé — Niveau Difficile}\\[2pt]
{\color{primary}\rule{0.5\textwidth}{1.2pt}}\\[4pt]
{\small\itshape Racine, quotient, mouvement et point anguleux : la définition jusqu'au bout}
\end{center}
\vspace{2mm}

\begin{corr}[1 --- $f(x)=\sqrt{x}$, $x_0>0$]
\textbf{a)} $\ds \tau(h)=\frac{\sqrt{x_0+h}-\sqrt{x_0}}{h}$.\\[3pt]
\textbf{b)} On multiplie numérateur et dénominateur par le conjugué
$\sqrt{x_0+h}+\sqrt{x_0}$ :
\[
\tau(h)=\frac{\big(\sqrt{x_0+h}-\sqrt{x_0}\big)\big(\sqrt{x_0+h}+\sqrt{x_0}\big)}
{h\big(\sqrt{x_0+h}+\sqrt{x_0}\big)}
=\frac{(x_0+h)-x_0}{h\big(\sqrt{x_0+h}+\sqrt{x_0}\big)}.
\]
Le numérateur se simplifie en $h$ :
\[
\tau(h)=\frac{h}{h\big(\sqrt{x_0+h}+\sqrt{x_0}\big)}
=\frac{1}{\sqrt{x_0+h}+\sqrt{x_0}}.
\]
\textbf{c)} On fait $h\to 0$ ; comme $\sqrt{x_0+h}\to\sqrt{x_0}$ :
\[
f'(x_0)=\frac{1}{\sqrt{x_0}+\sqrt{x_0}}=\boxed{\frac{1}{2\sqrt{x_0}}}.
\]
\textbf{d)} $\ds f'(4)=\frac{1}{2\sqrt{4}}=\frac{1}{2\times 2}=\boxed{\frac{1}{4}}$.
\end{corr}

\begin{corr}[2 --- mobile MRUA, $f(x)=\tfrac{1}{2}x^2+2x$]
\textbf{a)} $\ds \tau(h)=\frac{\big[\tfrac12(x_0+h)^2+2(x_0+h)\big]
-\big[\tfrac12 x_0^2+2x_0\big]}{h}$.\\[3pt]
\textbf{b)} On développe le numérateur :
\[
\tfrac12\big(x_0^2+2x_0h+h^2\big)+2x_0+2h-\tfrac12 x_0^2-2x_0
=x_0h+\tfrac12 h^2+2h=h\Big(x_0+\tfrac12 h+2\Big).
\]
Après simplification par $h$ : $\ds \tau(h)=x_0+\tfrac12 h+2$.\\[3pt]
\textbf{c)} En faisant $h\to 0$ : $\ds f'(x_0)=\boxed{x_0+2}$.\\[3pt]
\textbf{d)} À $t=3$~s : $f'(3)=3+2=\boxed{5\ \text{m/s}}$.
\end{corr}

\begin{corr}[3 --- $f(x)=\tfrac{1}{x+1}$, $x_0\neq -1$]
\textbf{a)} $\ds \tau(h)=\frac{1}{h}\left(\frac{1}{x_0+h+1}-\frac{1}{x_0+1}\right)$.\\[3pt]
\textbf{b)} On réduit au même dénominateur $(x_0+h+1)(x_0+1)$ :
\[
\frac{1}{x_0+h+1}-\frac{1}{x_0+1}
=\frac{(x_0+1)-(x_0+h+1)}{(x_0+h+1)(x_0+1)}
=\frac{-h}{(x_0+h+1)(x_0+1)}.
\]
Donc
\[
\tau(h)=\frac{1}{h}\cdot\frac{-h}{(x_0+h+1)(x_0+1)}
=\frac{-1}{(x_0+h+1)(x_0+1)}.
\]
\textbf{c)} En faisant $h\to 0$ (alors $x_0+h+1\to x_0+1$) :
\[
f'(x_0)=\frac{-1}{(x_0+1)(x_0+1)}=\boxed{-\frac{1}{(x_0+1)^2}}.
\]
\end{corr}

\begin{corr}[4 --- $f(x)=|x|$ en $x_0=0$]
Ici $f(0)=0$, donc $\ds \tau(h)=\frac{|0+h|-|0|}{h}=\frac{|h|}{h}$.\\[3pt]
\textbf{a)} Selon le signe de $h$ :
\begin{itemize}[leftmargin=6mm,itemsep=1pt]
  \item si $h>0$ : $|h|=h$, donc $\ds\tau(h)=\frac{h}{h}=1$ ;
  \item si $h<0$ : $|h|=-h$, donc $\ds\tau(h)=\frac{-h}{h}=-1$.
\end{itemize}
\textbf{b)} Les limites latérales sont donc
\[
\lim_{h\to 0^+}\tau(h)=1
\qquad\text{et}\qquad
\lim_{h\to 0^-}\tau(h)=-1.
\]
\textbf{c)} Les deux limites latérales existent mais sont \textbf{différentes}
($1\neq -1$) : la limite $\limh\tau(h)$ \emph{n'existe pas}. Par conséquent
$f$ \textbf{n'est pas dérivable en $0$}. Géométriquement, la courbe présente en $O$
un \textbf{point anguleux} : elle admet une demi-tangente de pente $1$ à droite et
une demi-tangente de pente $-1$ à gauche, sans tangente unique.

\begin{center}
\begin{tikzpicture}[scale=0.9,>=Latex]
  \draw[->,gray!70] (-2.4,0)--(2.4,0) node[right]{\small $x$};
  \draw[->,gray!70] (0,-0.4)--(0,2.3) node[above]{\small $y$};
  \draw[courbe,very thick] (-2,2)--(0,0)--(2,2);
  \fill[primarydark] (0,0) circle (1.8pt) node[below right,font=\scriptsize]{$O$};
  \node[courbe,font=\scriptsize] at (1.75,2.15) {$y=|x|$};
  \node[font=\scriptsize,tang] at (1.35,0.55) {pente $+1$};
  \node[font=\scriptsize,courbeder] at (-1.35,0.55) {pente $-1$};
\end{tikzpicture}
\end{center}
\end{corr}

\end{document}
